\section{Relation to the literature}\label{sec:literature}

This section says where each ingredient of the method comes from and what, as far as a targeted
but not exhaustive search shows, has not been done yet.

\subsection{Trajectory similarity with a temporal tolerance}

The LCSS of Vlachos, Kollios and Gunopulos~\cite{vlachos2002} matches two points when they are
closer than $\varepsilon$ in every coordinate and their \emph{indices} in the two sequences differ
by at most $\delta$. The similarity $S_1$ is the length of the common subsequence divided by the
length of the shorter track, and the variant $S_2$ maximises $S_1$ over all spatial translations,
to recognise parallel movements. The EDR of Chen, Özsu and Oria~\cite{chen2005} and DTW with a
Sakoe--Chiba band~\cite{sakoe1978} bound the shift in the same way. These measures return a
single number per pair of tracks: they say how similar two tracks are, and not \emph{when} two
pilots are together. A bound on indices counted from the first point of each track also mixes
the shift between the takeoffs with the delay in flight, and $S_2$ matches precisely the parallel
routes that we want to exclude.

The measures of joint movement that do say when are built on simultaneous fixes. The proximity
index Prox is the share of simultaneous fixes closer than a threshold~\cite[§3.1]{joo2018}: with
$\delta=0$ and without the conditions on height and heading, it is the share of the flight of $i$
in contact with $j$ in the sense of Section~\ref{sec:correspondences}.

\subsection{Following with a delay}\label{sec:following}

The closest precedent of our pair contact is the following of Li, Wu and Crofoot~\cite{li2013},
defined on tracks sampled at the same instants (Figure~\ref{fig:li}). At instant $i$ the local
lag $\Delta(i)$ is the shift $t$, with $|t|\le l_{\max}$, that minimises the distance between the
position $r_i$ of $R$ and the position $s_{i-t}$ of $S$. The instant is \emph{valid} if this
minimum distance is below $d_{\max}$, and $R$ follows $S$ at a valid instant if $\Delta(i)>0$
(§IV-A). The lag is therefore local and free to vary, and its sign separates the follower from
the leader. Two individuals that merely move together lead and follow equally often, so the
following intervals are the stretches with more following instants than the half expected
without following: the maximal segments of $f(i)-1/2$, with $f(i)=1$ when $R$ follows and $0$
otherwise, found in linear time~\cite{ruzzo1999} (§IV-B) and checked with the roles
reversed (§IV-C). The threshold $d_{\max}$ only discards the instants without interaction, and
does not move the minimiser.

Our contact keeps the window of lags and the distance threshold, and leaves the sign out: the
contact of $i$ with $j$ at instant $t$ is the validity of $t$, with tolerances on height and
heading added (Section~\ref{sec:correspondences}). We want the pairs that fly together, whether
or not one of them leads; who follows whom comes back afterwards, through the median lag of each
contact (Section~\ref{sec:phase3}).

\begin{figure}
\centering
\begin{tikzpicture}[x=1cm,y=1cm]
  \draw[track i, line width=8pt, blueI!22] plot[smooth] coordinates {(2.1,1.35) (3.4,1.85) (4.8,1.75) (6.2,2.25) (7.5,2.6)};
  \draw[track i] plot[smooth] coordinates {(0,0.6) (1.0,1.0) (2.1,1.35) (3.4,1.85) (4.8,1.75) (6.2,2.25) (7.5,2.6) (8.8,3.0) (10,3.2)};
  \foreach \p in {(1.0,1.0),(2.1,1.35),(3.4,1.85),(4.8,1.75),(6.2,2.25),(7.5,2.6),(8.8,3.0)} \fill[blueI] \p circle (1.3pt);
  \node[lbl, above left=-1pt] at (2.1,1.35) {$s_{i-l_{\max}}$};
  \node[lbl, above left] at (7.5,2.6) {$s_{i+l_{\max}}$};
  \node[pt i] (si) at (4.8,1.75) {};
  \node[lbl, blueI, below right=-1pt] at (si) {$s_i$};
  \draw[track j] plot[smooth] coordinates {(0,0.0) (1.2,0.45) (2.5,0.95) (3.35,1.35)};
  \node[pt j] (ri) at (3.35,1.35) {};
  \node[lbl, orangeJ, below right] at (ri) {$r_i$};
  \node[pt i] (sd) at (3.4,1.85) {};
  \draw[ink, densely dashed] (ri) -- (sd);
  \node[lbl, above left] at (sd) {$s_{i-\Delta(i)}$};
  \node[lbl, blueI] at (10.3,3.2) {$S$};
  \node[lbl, orangeJ] at (-0.3,0.0) {$R$};
  \draw[decorate, decoration={brace, amplitude=4pt}, ink2] (2.1,2.35) -- (7.3,3.25)
    node[midway, above=6pt, lbl, ink, sloped] {points of $S$ recorded from $i-l_{\max}$ to $i+l_{\max}$};
  \node[note, align=left, anchor=west] at (8.4,1.2) {closest point: $s_{i-\Delta(i)}$\\with $\Delta(i)>0$: $S$ passed\\there first, so $R$ follows $S$};
\end{tikzpicture}
\caption{Following with a local lag~\cite{li2013}. At instant $i$, the track of $S$ is searched
in the time window $[i-l_{\max}, i+l_{\max}]$ (shaded) for the point closest to the position
$r_i$ of $R$. The best lag $\Delta(i)$ is positive: $R$ is now where $S$ was $\Delta(i)$
earlier.}
\label{fig:li}
\end{figure}

Nagy et al.~\cite{nagy2010} measure the delay between two pigeons as the shift that maximises the
correlation of their headings, and build leader--follower hierarchies from its sign. In a thermal
the heading turns through $2\pi$ every 20--30~s, so this correlation, like the distance of
Section~\ref{sec:roles}, repeats with the period of the turn. Konzack et al.~\cite{konzack2017}
display such local delays between simultaneously recorded tracks in a ``delay space'', comparing
Fréchet, DTW and edit-distance matchings. FLICA and its package
mFLICA~\cite{amornbunchornvej2018a,amornbunchornvej2018b,amornbunchornvej2021} infer a following
network in sliding time windows, with DTW restricted to a Sakoe--Chiba band, and call the faction
of a leader the individuals that reach it through a directed path: groups built from delayed
following, but directed, with leaders, and on windows. The following of Andersson et
al.~\cite{andersson2007} is instantaneous: an individual follows another that is in a sector
ahead of it with a similar heading.

\subsection{Group structure in time}

Buchin et al.~\cite{buchin2015} call two entities directly connected at time $t$ if their
$\varepsilon$-discs overlap, and $\varepsilon$-connected if a chain of direct connections joins
them at that instant; the maximal $\varepsilon$-connected sets, the components, partition the
entities at every instant. The Reeb graph records how the components evolve: its vertices are the
events (start, end, merge, split), and each edge is a fixed set of entities that forms one
component over an interval of time (Figure~\ref{fig:reeb}). On it they define the maximal groups,
sets of at least $m$ entities that stay inside one component for at least a duration $\delta$;
maximal groups may overlap, and an entity can belong to several at once. Two more ideas of theirs
reappear below: a robust structure ignores interruptions shorter than $\alpha$, by calling two
entities connected at $t$ when they were directly connected at some instant within $\alpha/2$ of
$t$ (§4); and the grouping structure can split the trajectories into independent data sets
(§1). The proximity is always simultaneous.

\begin{figure}
\centering
\begin{tikzpicture}[x=1cm,y=1cm,
    vtx/.style={circle, fill=ink, inner sep=1.5pt},
    single/.style={muted, line width=0.9pt},
    multi/.style={blueI, line width=2.6pt},
    set/.style={font=\footnotesize, fill=white, inner sep=1pt}]
  \coordinate (s1) at (0,2.2); \coordinate (s2) at (0,1.0); \coordinate (s3) at (0,0.0);
  \coordinate (m1) at (3.0,1.6); \coordinate (sp) at (6.0,1.6); \coordinate (m2) at (9.0,0.5);
  \coordinate (e1) at (12.5,2.3); \coordinate (e2) at (12.5,0.5);
  \draw[multi] (s1) to[out=0, in=170] node[set, pos=0.45, above=1pt] {$\{a,b\}$} (m1);
  \draw[single] (s2) to[out=0, in=190] node[set, pos=0.45, below=1pt] {$\{c\}$} (m1);
  \draw[multi] (m1) -- node[set, above=1pt] {$\{a,b,c\}$} (sp);
  \draw[single] (sp) to[out=10, in=180] node[set, pos=0.55, above=1pt] {$\{a\}$} (e1);
  \draw[multi] (sp) to[out=-10, in=170] node[set, pos=0.45, above right=0pt] {$\{b,c\}$} (m2);
  \draw[single] (s3) to[out=0, in=190] node[set, pos=0.5, below=1pt] {$\{d\}$} (m2);
  \draw[multi] (m2) -- node[set, above=1pt] {$\{b,c,d\}$} (e2);
  \foreach \v in {s1,s2,s3,m1,sp,m2,e1,e2} \node[vtx] at (\v) {};
  \node[note, below=3pt] at (m1) {merge};
  \node[note, below=3pt] at (sp) {split};
  \node[note, below=3pt] at (m2) {merge};
  \draw[->, ink2] (-0.2,-0.8) -- (12.9,-0.8) node[right, lbl, ink] {$t$};
\end{tikzpicture}
\caption{The grouping structure of Buchin et al.~\cite{buchin2015} as a Reeb graph. Each edge is a
fixed set of entities that forms one connected component over an interval; vertices are starts,
ends, merges and splits. The thick edges have at least two entities: in this note they become the
configurations (Section~\ref{sec:configurations}), built on delayed contacts.}
\label{fig:reeb}
\end{figure}

\subsection{Interaction indices and null models}

Joo et al.~\cite{joo2018} sort the metrics of joint movement between pairs by what they capture,
proximity or coordination. Long et al.~\cite{long2014} compare eight of them on simulated tracks
and on deer: the indices that rest on a statistical test report interactions that are not there,
more often the finer the sampling, and the indices that return one number per pair miss the
interactions that come and go; the wildlifeDI package~\cite{long2022} extends the analysis to
contacts in high-frequency data. Spiegel et al.~\cite{spiegel2016} build a null model by shuffling
whole days of each individual's path among its days: each individual keeps its use of space, and
the synchrony between individuals is lost.

\subsection{Collective soaring}

Ákos, Nagy and Vicsek~\cite{akos2008} compare the soaring of birds with that of paraglider and
hang-glider pilots on GPS tracks; falcons, storks and paragliders circle with radii of 20--26~m.
In migrating white storks a few leaders find and explore the thermals and the others follow them,
leaving the thermals earlier and lower~\cite{flack2018}; inside a thermal, flock mates
synchronise their circling~\cite{nagy2018}. Vultures informed by other birds of the next updraft
glide towards it faster~\cite{williams2018}. We found no work that identifies group flights
systematically in large free-flight datasets.

\subsection{What we take and what is new}

The method is not a new general algorithm for grouping trajectories. The pair contact is the
proximity with a local lag of Li et al.~\cite{li2013}, the group structure is that of Buchin et
al.~\cite{buchin2015}, and groups built from delayed relations exist in
mFLICA~\cite{amornbunchornvej2021}, with other definitions (Table~\ref{tab:origin}). What we did
not find in the literature is:
\begin{enumerate}[label=(\roman*)]
  \item the Reeb structure built on contacts with a local delay instead of simultaneous proximity;
  \item the adaptations to thermal flight: a vertical tolerance that separates pilots circling in
        the same thermal at distant heights, and headings compared only in glide
        (Section~\ref{sec:roles});
  \item the application to a free-flight dataset of this size, with thresholds calibrated on a
        null model.
\end{enumerate}
The contribution is therefore a physically motivated adaptation and a scalable implementation,
more than a new method in itself, and the thesis should present it in these terms.

\begin{table}
\centering
\small
\begin{tabularx}{\textwidth}{@{}>{\raggedright}p{2.1cm} >{\raggedright}p{2.6cm} >{\raggedright}X >{\raggedright\arraybackslash}X@{}}
\toprule
Component & Precedent & What we take & What we change \\
\midrule
Pair contact with a delay & Li et al.~\cite{li2013}; tolerance in~\cite{vlachos2002,chen2005,sakoe1978} &
  window of lags and distance threshold; local lag, free to vary &
  height tolerance (separates heights within a thermal); heading only in glide; local phase
  excluded; no sign in the contact, $a_{ij}=a_{ji}$ (the sign returns in $\ell_{\rm med}$) \\ \addlinespace
Contact intervals & robustness in~\cite{buchin2015}; maximal segments in~\cite{li2013} &
  tolerance to short interruptions &
  gaps shorter than $g$ filled; minimum duration calibrated on a null model \\ \addlinespace
Groups in time & Buchin et al.~\cite{buchin2015} & components and their Reeb graph; configurations
  are its edges with at least two entities &
  built on delayed contacts; episodes are the connected parts of the Reeb graph without single
  entities; no maximal groups \\ \addlinespace
Groups from following & mFLICA~\cite{amornbunchornvej2021} & -- &
  undirected groups without leaders; exact intervals instead of sliding windows \\ \addlinespace
Null model & Spiegel et al.~\cite{spiegel2016}; SAM~\cite{tusher2001} & time shift of the tracks;
  false discovery rate as chance count over observed count &
  shift within the day, drawn from the real delays between takeoffs; chance count per unit of
  exposure \\ \addlinespace
Computation & filter--refine scheme & -- &
  no-loss lemma for delayed contacts; KD-tree with a time coordinate \\
\bottomrule
\end{tabularx}
\caption{Where the ingredients of the method come from.}
\label{tab:origin}
\end{table}

\paragraph{Essential reading.} The two readings to do first are Li et al.~\cite{li2013},
Sections~II and~IV (following, local lag, score), and Buchin et al.~\cite{buchin2015}, Sections~1,
2.1 and~4 of the arXiv version (groups, Reeb graph, robustness). Long et al.~\cite{long2014} and
Spiegel et al.~\cite{spiegel2016} serve the null model; Vlachos et al.~\cite{vlachos2002} and Joo
et al.~\cite{joo2018} place the measures of similarity and proximity.
